% ---------------------------------------------------------------------
%  第 10 章 平衡态与稳态
% ---------------------------------------------------------------------
\chapter{平衡态与稳态}
\label{ch:equilibrium}

第\ref{ch:solvers}章假定初始 Green 函数已知。本章说明这些初值
如何由物理输入构造：哈密顿量、态密度、单极点谱函数，以及它们在
稳态形式下的对应物。所有构造都是解析的，不需要迭代；它们的精度
只受数值积分方法限制。

\section{哈密顿量构造}
\label{sec:from-ham}

最直接的构造从自由哈密顿量 $\bm{H}_0$ 出发~\cite{aoki2014}。先对角化
\begin{equation}
    \bm{H}_0 = \bm{R}\,\mathrm{diag}(\lambda_1,\dots,\lambda_d)\,
    \bm{R}^{\dagger},
\end{equation}
再按能级与费米/玻色分布写出四个分量。对费米型，
\begin{align}
    \Gmat(\tau) &= -\bm{R}\,
        \mathrm{diag}\big[\ee^{\tau\lambda_i}n_F(\lambda_i)\big]\,
        \bm{R}^{\dagger},\\
    \Glmix(t_n,\tau) &= \ii\,\bm{U}_n\bm{R}\,
        \mathrm{diag}\big[\ee^{-\tau\lambda_i}n_F(-\lambda_i)\big]\,
        \bm{R}^{\dagger},\\
    \Gret(t_n,t_j) &= -\ii\,\bm{U}_n\bm{U}^{\dagger}_j,\\
    \Gless(t_j,t_n) &= \ii\,\bm{U}_j\,\bm{N}\,\bm{U}^{\dagger}_n,
    \qquad \bm{N} = \bm{R}\,
        \mathrm{diag}\big[n_F(-\lambda_i)\big]\bm{R}^{\dagger}.
    \label{eq:from-ham}
\end{align}
对玻色型，$\Gmat$ 与 $\bm{N}$ 变号，$n_F$ 换成 $n_B$，
$\Glmix$ 乘 $-1$。

\paragraph{传播子。}
常数哈密顿量时 $\bm{U}_n=\ee^{\ii hn\bm{K}}$，其中
$\bm{K}=\mu-\bm{H}_0$。含时哈密顿量时使用无对易子展开：
\begin{align}
    \bm{U}_n &= \ee^{-\ii h\Hm(t_{n-1/2})}\,\bm{U}_{n-1},
        && \text{二阶},\\
    \bm{U}_n &= \ee^{-\ii h(a_1\Hm_1+a_2\Hm_2)}
        \ee^{-\ii h(a_2\Hm_1+a_1\Hm_2)}\,\bm{U}_{n-1},
        && \text{四阶},
\end{align}
其中 $c_{1,2}=\tfrac12\mp\tfrac{\sqrt3}{6}$，
$a_{1,2}=(3\mp2\sqrt3)/12$。中间时刻的哈密顿量由 $k$ 阶多项式
插值给出。最后把 $\bm{U}_n$ 乘上 $\ee^{\ii\mu t_n}$。

\section{态密度构造}
\label{sec:from-dos}

当体系由态密度 $A(\epsilon)$ 描述时，Green 函数的各分量是能量积分：
\begin{align}
    \Gret(t) &= -\ii\int\!\dd\epsilon\,A(\epsilon)\,
        \ee^{-\ii(\epsilon-\mu)t},\\
    \Gless(t) &= -\xi\,\ii\int\!\dd\epsilon\,A(\epsilon)\,
        n_\xi(\epsilon-\mu)\,\ee^{-\ii(\epsilon-\mu)t},
\end{align}
以及虚时间上的对应表达式。实现用自适应 Fourier 积分：积分区间
按“误差估计最大的面板优先细分”构造，面板数上限由 \code{limit}
给出，每个面板含 \code{nn} 个子区间。

\paragraph{参数约束。}
\code{nn} 必须为偶数且不小于 $20$；\code{limit} 至少为 $1$。
态密度在每个采样点必须有限且非负；玻色型在 $\epsilon=\mu$ 处的
零模被显式拒绝，因为 $n_B$ 在那里发散。

\paragraph{Bethe 晶格。}
半圆态密度
\begin{equation}
    A(\epsilon) = \frac{2}{\pi W^2}
    \sqrt{W^2-(\epsilon-\epsilon_0)^2},
    \qquad |\epsilon-\epsilon_0|\le W
\end{equation}
由 \code{BetheDos} 提供，常用于检验动力学平均场理论
（Dynamical Mean-Field Theory，DMFT）的实现。它的归一化
$\int A=1$ 意味着等时推迟分量满足 $\Gret(t,t)=-\ii$，这一性质
在测试中被显式检查。

\section{单极点构造}
\label{sec:pole}

单极点 Green 函数描述一个孤立的玻色模式：
\begin{equation}
    \Gret(\omega) = \frac12\left[
        \frac{1}{\omega-w+\ii 0^+} - \frac{1}{\omega+w+\ii 0^+}\right].
\end{equation}
令 $u=t-t'\ge0$。库保存的小于块与推迟块分别为
\begin{align}
    \Gless(t',t) &= -\tfrac{\ii}{2}\Big[
        (1+n_B(w))\ee^{-\ii wu}+n_B(w)\ee^{\ii wu}\Big],\\
    \Gret(t,t') &= -\sin(wu).
\end{align}
实现要求 $w\bet>0$（否则玻色分布无定义）且对象为玻色型。该构造
常用于 Holstein 模型的声子自由度~\cite{werner2009}。

\section{Bethe 晶格}
\label{sec:bethe}

Bethe 晶格的半圆态密度给出解析的自由传播子~\cite{freericks2006}。在 DMFT 中，它把
晶格问题化为一个有效的单杂质问题：杂质自能与局域 Green 函数之间
由自洽条件联系。库提供该构造作为算例与测试的共同基础，
\cref{fig:dmft-loop} 给出 DMFT 的自洽环。

\begin{figure}[htbp]
    \centering
    \begin{tikzpicture}[
        node distance=10mm,
        box/.style={draw=kbblue!60, fill=kbblue!6, rounded corners=3pt,
                    font=\small, align=center, inner sep=5pt,
                    minimum height=9mm, text width=3.1cm},
        ar/.style={-{Stealth[length=6pt]}, thick, kbblue!70}
    ]
        \node[box] (g) {局域 Green 函数\\$\Gless(t,t')$};
        \node[box, right=of g] (sig) {自能\\$\Sig$};
        \node[box, below=of sig] (dyson) {Dyson 方程\\$[\Gret_0]^{-1}-\Sig$ 的逆};
        \node[box, left=of dyson] (dos) {态密度/Bethe\\$\Gret_0$};

        \draw[ar] (g) -- (sig);
        \draw[ar] (sig) -- (dyson);
        \draw[ar] (dyson) -- (dos);
        \draw[ar] (dos) -- (g);
    \end{tikzpicture}
    \caption[DMFT 自洽环]{DMFT 自洽环。库提供环上的每个构件；算例把环组织成
             迭代。}
    \label{fig:dmft-loop}
\end{figure}

\section{稳态 Dyson 方程}
\label{sec:steady-dyson}

稳态推迟分量由频域 Dyson 方程给出：
\begin{equation}
    \Gret(\omega) =
    \Big[\omega+\mu-\bm{H}-\Sigret(\omega)
      -\Sigret_{b}(\omega)\Big]^{-1},
\end{equation}
小于分量由
\begin{equation}
    \Gless(\omega) =
    \Gret(\omega)\big[\Sigless(\omega)+\Sigless_b(\omega)\big]
    \Gret(\omega)^{\dagger}
\end{equation}
给出。实现先把自能变换到频域（在局部副本上做，不修改输入），
填充中心频带，再变换回时间域。中心频带的宽度为
\begin{equation}
    |m| < \frac{N_{ft}/2-1}{3},
\end{equation}
其余模式置零。这一截断是稳态计算精度的主要来源。

\section{浴自能}
\label{sec:bath}

\cref{tab:bath}列出各类浴模型及其自能形式。

浴自能有三种解析形式，由 \code{DysonBathType} 选择：
\begin{table}[htbp]
    \centering
    \caption[浴自能]{浴自能}
    \label{tab:bath}
    \small
    \begin{tabularx}{\textwidth}{@{}llX@{}}
        \toprule
        类型 & 参数 & 形式 \\
        \midrule
        \code{Const} & $\eta$ &
            $\Sigret=-\ii\eta$，$\Sigless=2\ii\eta\,n_F(\omega-\mu)$ \\
        \code{Gauss} & $\eta$ &
            $\Sigret=-\ii\eta\ee^{-\omega^2}$，$\Sigless=2\ii\eta\ee^{-\omega^2}n_F$ \\
        \code{Ohmic} & $\eta$ &
            $\Sigret=-\ii\eta\,\omega\ee^{-\omega^2}$，
            $\Sigless$ 由 $n_B$ 与 $\omega$ 组合给出 \\
        \bottomrule
    \end{tabularx}
\end{table}

高斯与欧姆型都带有 $\ee^{-\omega^2}$ 的阻尼因子。代码对其形式做了
逐字保留：指数为 $-\omega^2$，与时间步长无关；\code{h} 只做合法性
检查，不参与计算。欧姆型在 $\omega\to0$ 处用极限值
\begin{equation}
    \Sigless \to -\frac{2\ii\eta}{\bet}\ee^{-\omega^2}
\end{equation}
避免 $1/\sinh(\bet\omega)$ 的相消。

\section{态密度计算}
\label{sec:steady-dos}

由稳态 Green 函数提取态密度有两条路径：
\begin{itemize}
    \item 频域直接读取：$\mathrm{Im}\Gret(\omega)=-\pi A(\omega)$，
          适用于 \code{Trapez} 变换后的数据；
    \item 自适应积分：给定 $A(\omega)$，用自适应 Fourier 积分直接
          构造时间域的 $\Gret$ 与 $\Gless$。
\end{itemize}
两条路径互为逆运算，测试中用于检查一致性。

\section{平衡化}
\label{sec:equilibrate}

\code{equilibrate_lesser} 把小于分量替换为满足涨落耗散关系的平衡值。
给定推迟分量与化学势，平衡小于分量为
\begin{equation}
    \Gless(\omega) = \xi\,
    \big[\Gret(\omega)-\Gret(\omega)^{\dagger}\big]\,
    n_\xi(\omega-\mu),
\end{equation}
其中 $\xi=-1$ 对应费米型。实现按模式逐个计算，Nyquist 点置零。
该函数用于把任意初值投影到平衡态，也用于构造参考解。

\section{稳态密度矩阵与气泡图}
\label{sec:steady-obs}

稳态密度矩阵由等时小于分量给出（$\xi$ 为统计符号）：
\begin{equation}
    \rho = \ii\,\xi\,\Gless(t=0),
\end{equation}
接口直接读取时域的零相对时间块。对物理归一化的频域数据，梯形逆变换
在零时刻的求和因子为 $\Delta\omega/(2\pi)=1/(N_{ft}h)$。

气泡图也有稳态版本：
\code{bubble1_steady} 与 \code{bubble2_steady}
在频域中按点乘构造极化与自能，随后用逆变换回到时间域。
标量变体用于 $1\times1$ 轨道，以省去块循环。

\section{网格工具}
\label{sec:grid-tools}

\code{SteadyGrid} 提供三个工具：
\begin{itemize}
    \item \code{time_at(m)} 与 \code{freq_at(m)}：把模式索引换算成
          时间或频率，负模式按周期回绕；
    \item \code{set_h(h)}：修改时间步并自动重算
          $\Delta\omega=\pi/((N_{ft}/2)h)$；
    \item \code{set_dw(dw)}：修改频率步并自动重算时间步。
\end{itemize}
两个步长互为倒数关系，任何修改都保持网格自洽。Nyquist 模式
（$m=N_{ft}/2$）是唯一自共轭的模式，所有变换与对称性操作都把它
单独处理。

\section{由 Matsubara 分量启动}
\label{sec:from-mat}

在时间推进开始之前，实时间分量由 Matsubara 分量播种：
\begin{align}
    \Glmix(0,\tau_m) &= \ii\,\xi\,\Gmat(\bet-\tau_m),\\
    \Gless(0,0) &= \Glmix(0,0),\\
    \Gret(0,0) &= \ii\,\Gmat(0)-\Gless(0,0).
\end{align}
\code{from_mat} 填充第 $0$ 步，\code{from_mat_k} 填充前 $k$ 步。
后者用于启动阶段：当求解器需要前 $k$ 个时间步的初值时，它比逐个
调用 \code{from_mat} 更高效，也保证边界值一致。
